% Ad Familiares 14.1 --- headnote
% ad-familiares-14-01, 25 November 58 BC, written at Dyrrachium

\textit{Cicero to Terentia, Tullia, and the boy Marcus, written
from Dyrrachium on the sixth day before the Kalends of
December (25 November) 58 BC, in the same week as the move
from Thessalonica that the parallel \textit{Att.} 3.22
records. The crisis at Rome is the new tribunes' bill: by
December the bid for Cicero's recall would be tested. The
opening note is what reaches Cicero of Terentia's bearing
``incredible courage and strength of spirit,'' and the
self-blame --- not fate, but his own fault, in trusting men
who envied him.}

\textit{The painful detail in \textsection 5 is the news that
Terentia means to sell their \textit{vicus} (a village
property, possibly from her own \textit{dos}) to fund the
recall effort: Cicero begs her not to. ``If our friends shall
stand to their duty, money will not be lacking; if they shall
not, you will not be able to make it good with your money.''
The closing thought is for the future of young Marcus: ``see
that we do not destroy the boy in his ruin.''}

\textit{The closing \textsection 7 explains the move:
Dyrrachium is a free city, friendly to him, and the nearest
landing on the Adriatic to Italy --- proximity is what now
matters. The exile letters thin out from here: Cicero will
spend the early winter at Dyrrachium and then begin moving
again as the senatorial decrees of 1 January 57 BC and Lentulus
Spinther's consulship open the road home.}
